% ============================================================================
%  Sovereign Bunny Stack -- Packaging, Runtime and Deployment (component paper)
%  SKELETON: outline and content-migration map only. Base: sovereign-booklet.tex
% ============================================================================
\documentclass[a4paper,11pt,twoside,parskip=half-,headings=small]{scrartcl}
% ============================================================================
%  common-preamble.tex -- shared packages, colours, boxes and status macros for
%  the Sovereign Bunny paper set (all papers except the base booklet, which
%  keeps its own preamble). Each paper sets its own \documentclass, hyperref
%  metadata and fancyhdr heads (see \setpaperheads).
% ============================================================================
\usepackage[T1]{fontenc}
\usepackage[utf8]{inputenc}
\usepackage{lmodern}
\usepackage[english]{babel}
\usepackage{microtype}
\setlength{\emergencystretch}{2em}
\usepackage[margin=2.5cm,headheight=14pt]{geometry}
\usepackage{amsmath,amssymb}
\usepackage{booktabs,tabularx,longtable,array}
\usepackage{xcolor}
\usepackage{tikz}
\usetikzlibrary{arrows.meta,positioning,shapes.geometric,fit,calc,backgrounds}
\usepackage{caption}
\usepackage{needspace}
\usepackage{enumitem}
\usepackage{tcolorbox}
\tcbuselibrary{breakable}
\usepackage{listings}
\usepackage{fancyhdr}
\usepackage{float}
\usepackage{ragged2e}
\usepackage[hidelinks]{hyperref}
\usepackage[capitalise,noabbrev]{cleveref}

\setlength{\parindent}{0pt}
\captionsetup{font=small,labelfont=bf,justification=centering}

\definecolor{ink}{HTML}{1B2A41}
\definecolor{accent}{HTML}{0B7A75}
\definecolor{warn}{HTML}{B5472B}
\definecolor{mute}{HTML}{5C6B7A}
\definecolor{codebg}{HTML}{F8FAFC}
\addtokomafont{section}{\color{ink}}
\addtokomafont{subsection}{\color{ink}}
\addtokomafont{subsubsection}{\color{ink}}
\setlist{itemsep=2pt,topsep=4pt}

\lstdefinestyle{customcode}{
  backgroundcolor=\color{codebg},commentstyle=\color{mute}\itshape,
  keywordstyle=\color{accent}\bfseries,stringstyle=\color{warn},
  basicstyle=\ttfamily\scriptsize,breaklines=true,captionpos=b,keepspaces=true,
  showstringspaces=false,frame=single,rulecolor=\color{mute!30},tabsize=2}
\lstdefinelanguage{Rust}{
  keywords={true,false,struct,enum,fn,let,mut,pub,impl,for,if,else,match,return,const,u8,u16,u32,u64,usize,bool,String,Vec,Result,Ok,Err,Option,Some,None,use,self,Self},
  sensitive=true,comment=[l]{//},morecomment=[s]{/*}{*/},morestring=[b]"}
\lstdefinelanguage{JSON}{
  string=[s]{"}{"},stringstyle=\color{warn},comment=[l]{//},
  morekeywords={true,false,null},keywordstyle=\color{accent}\bfseries}
\lstdefinelanguage{TypeScript}{
  keywords={export,interface,type,const,let,function,async,await,return,import,from,new,readonly,number,string,bigint,boolean,void,Promise,Map,Uint8Array,ReadableStream,true,false},
  sensitive=true,comment=[l]{//},morecomment=[s]{/*}{*/},morestring=[b]",morestring=[b]'}
\lstset{style=customcode}

\pagestyle{fancy}
\fancyhf{}
\fancyfoot[C]{\small\thepage}
\renewcommand{\headrulewidth}{0.3pt}

\IfFileExists{build/meta.tex}{\input{build/meta.tex}}{%
  \newcommand{\bookletcommit}{unversioned}%
  \newcommand{\bookletdate}{\today}}

\newcommand{\EUR}[1]{\texteuro\,#1}
\newcommand{\code}[1]{\texttt{#1}}
\newcommand{\claim}[1]{\textcolor{warn}{[#1]}}
\newcommand{\stImpl}{\textcolor{accent}{\textbf{Implemented}}}
\newcommand{\stPart}{\textcolor{mute}{\textbf{Partial}}}
\newcommand{\stStub}{\textcolor{warn}{\textbf{Stub / unsafe}}}
\newcommand{\stSpec}{\textcolor{mute}{\textbf{Specified}}}
\newcommand{\stPlan}{\textcolor{mute}{\textbf{Planned}}}
\newcommand{\stRes}{\textcolor{warn}{\textbf{Research}}}
\newcolumntype{L}[1]{>{\RaggedRight\arraybackslash}p{#1}}
\newcolumntype{Y}{>{\RaggedRight\arraybackslash}X}

\newtcolorbox{keybox}[2][]{%
  breakable,colback=accent!6,colframe=accent,boxrule=0.6pt,arc=2pt,
  left=7pt,right=7pt,top=3pt,bottom=3pt,
  fonttitle=\bfseries,coltitle=accent,colbacktitle=accent!6,
  title={#2},before skip=4pt,after skip=4pt,#1}
\newtcolorbox{warnbox}[2][]{%
  breakable,colback=warn!5,colframe=warn,boxrule=0.6pt,arc=2pt,
  left=7pt,right=7pt,top=4pt,bottom=4pt,
  fonttitle=\bfseries,coltitle=warn,colbacktitle=warn!5,
  title={#2},before skip=6pt,after skip=6pt,#1}

\raggedbottom
\newcommand{\setpaperheads}[2]{\fancyhead[LE,RO]{\small\color{mute}#1}\fancyhead[RE,LO]{\small\color{mute}#2}}
\newcommand{\booklet}{the Sovereign Bunny whitebooklet}

\hypersetup{pdftitle={Sovereign Bunny Stack: Packaging, Runtime and Deployment},pdfauthor={Sovereign Bunny contributors},pdfsubject={Hermetic builds, SGX enclaves, deployment, end-to-end types},pdfkeywords={sovereign bunny, stack, nixos, sgx, gramine, argocd, sbom}}
\setpaperheads{Sovereign Bunny Stack}{Sovereign Bunny paper set -- Stack}

\begin{document}
\begin{titlepage}
  \centering
  \vspace*{2.2cm}
  {\color{mute}\scshape\large Component Paper (skeleton)\par}
  \vspace{0.6cm}
  {\color{ink}\Huge\bfseries Sovereign Bunny Stack\par}\vspace{0.4cm}
  {\color{accent}\Large Packaging, Runtime and Deployment\par}
  \vspace{0.8cm}
  {\large Hermetic builds, measured enclaves, supply-chain gates, and the operator's sovereign instance\par}
  \vfill
  \begin{tabular}{rl}
    \textbf{Status}   & Skeleton -- content still lives in the whitebooklet\\
    \textbf{Base}     & Sovereign Bunny whitebooklet v0.5\\
    \textbf{Revision} & \code{\bookletcommit}\\
    \textbf{Built}    & \bookletdate\\
  \end{tabular}
  \vspace*{2.2cm}
\end{titlepage}

\section{Purpose}
The Stack is how the ledger is \emph{operated}: reproducible builds, signed and measured artefacts, deployment pipelines, and schemas carried unchanged from contracts to the browser.
It makes no claims about consensus or economics. Its central contract is the \textbf{exit drill}: any operator can export state and restore on bare metal with bit-for-bit state-root equality, continuously, in CI.

\section{Proposed outline and migration map}
\begin{table}[H]
\centering\small
\begin{tabularx}{\textwidth}{@{}L{4.6cm}Y@{}}
\toprule
\textbf{Section of this paper} & \textbf{Source in the whitebooklet (by topic)} \\
\midrule
Cluster versus Stack, and the local sovereign instance & Cluster vs.\ Stack subsection \\
Hermetic NixOS builds and ArgoCD delivery & NixOS packaging, ArgoCD, end-to-end types \\
Measured enclaves: what TEEs can and cannot do & Hardware realism; Gramine manifest \\
Minimal Viable Security Kernel & MVSK section \\
CRA / NIS2 evidence and SBOM witnesses & CRA and NIS2; CycloneDX SBOM parser and witness generation \\
Exit drill and operational tests & Operational principles and tests (exit, reproducibility) \\
Zero-open-port operator access (zkSSH, BGP peering) & Sovereign SSH; BGP peering sections \\
Honest status of the reference implementation & Status register; regression deconstruction \\
Reference client (local-first SDK) & New: see \cref{sec:ref_client}; builds on the end-to-end type-preservation material \\
\bottomrule
\end{tabularx}
\caption{Content-migration map for the Stack paper.}
\end{table}

\section{Reference client implementation}
\label{sec:ref_client}

The ledger paper fixes four framework-neutral client invariants. This section describes how the Stack realises them today and what is planned. The TypeScript/WASM and Lit/XState choices below are \emph{reference}, not normative: another host that satisfies the invariants is equally conformant. This completes the end-to-end type-preservation pipeline (Rust and Solidity, through WASM bindings, to TypeScript) with a local-first state layer.

\subsection{Four layers, one responsibility each}
React Server Components blurred separation of concerns by mixing network calls and database queries into the view layer. The client stack avoids that with a strict vertical split.

\begin{table}[H]
\centering\small
\begin{tabularx}{\textwidth}{@{}L{3.3cm}YL{2.6cm}@{}}
\toprule
\textbf{Layer} & \textbf{Responsibility and invariant} & \textbf{Status} \\
\midrule
1. Local state and persistence & Application code mutates a typed actor store and configures a storage adapter (IndexedDB or Origin Private File System in a browser; SQLite or an embedded key-value store on a node). It contains no RPC calls, no circuit imports and no consensus logic, and works fully offline & \stPlan{} (wallet state machines exist: \path{wallet/ts/src/lattice_machine.ts}, \path{storage_manager.ts}) \\
2. Reconciler (WASM, worker) & Maintains the double-buffered Virtual CAR; on commit diffs against the last settled roots, updates accumulators, computes $\Delta_{\text{slot}}$ and the next nonce & \stPlan{} (Hive P2) \\
3. Prover and signer bridge & Selects the circuit for each dirty slot (\path{wallet/ts/src/circuits.ts}), formats witnesses, proves, and requests one signature from a hardware or wallet signer & \stPart{} (circuit loading and wallet connector exist) \claim{verify scope} \\
4. Transport and sync adapter & Packs the sealed frame as a \code{BNY\textbackslash x01} SSZ envelope and dispatches it onto the urgent or background lane without blocking local reads & \stPart{} (envelope tooling in \path{crates/cli}, \path{wallet/ts/src/envelopes.ts}) \claim{verify} \\
\bottomrule
\end{tabularx}
\caption{Reference client layers. Status words follow the whitebooklet; scopes marked verify have not been audited for this draft.}
\end{table}

\subsection{Interface sketch}
\begin{lstlisting}[language=TypeScript,caption={Specification sketch: host-neutral client interface (target file \code{wallet/ts/src/car\_reconciler.ts}; not yet present).},label={lst:car_client}]
// Storage and transport are pluggable host capabilities.
export interface LocalStoreAdapter {
  readSlot(slot: number): Promise<Uint8Array>;
  commitAtomicBatch(nonce: bigint, dirty: Map<number, Uint8Array>): Promise<void>;
}
export interface SyncTransportAdapter {
  dispatchFrame(frame: Uint8Array, lane: "urgent" | "background"): Promise<EpochReceipt>;
}

// Mirrors the normative lifecycle; Rejected is reachable from Staged and Sealed.
export type CommitPhase =
  | { phase: "Staged";   nonce: bigint }
  | { phase: "Sealed";   nonce: bigint; frame: Uint8Array }
  | { phase: "Anchored"; nonce: bigint; epoch: bigint }
  | { phase: "Rejected"; nonce: bigint; reason: string; rolledBackFrom: bigint };

export interface SovereignActorClient<TSchema> {
  // Owned, single-writer slots only.
  reconcile(mutator: (draft: TSchema) => void,
            onPhase?: (s: CommitPhase) => void): Promise<CommitPhase>;
  // Contended, cross-actor: always an explicit state machine, never a bare promise.
  proposeIntent(intent: CrossActorIntent): ReadableStream<SagaState>;
}
\end{lstlisting}
Package names such as a local-state library or an Iroh sync adapter are proposals; none of them is published or exists in the workspace today.

\subsection{Three edge cases the SDK must not hide}
\begin{enumerate}
  \item \textbf{Owned state versus multi-party contention.} Owned-slot writes commit locally and deterministically because nobody else can reject them (subject to the writer-lease rule). Anything touching another actor can fail on lock contention. The API therefore separates \code{reconcile} from \code{proposeIntent}, and the latter returns a state machine, so a developer cannot confuse a local write with a global lock.
  \item \textbf{The latency gap between a JavaScript write and a proven, anchored one.} A local write is microseconds to milliseconds; browser proving and epoch finality are orders of magnitude slower (proving time and the ten-second epoch are targets, not measurements, for this stack \claim{unmeasured}). The status of every state atom is therefore reactive and exposes Staged, Sealed, Anchored and Rejected. Staged data must be durable before the call returns, and a resumed session must be able to find Staged work after a tab closes and finish sealing it, so closing a tab never silently loses an accepted edit.
  \item \textbf{Host-language numerics and ordering.} Canonical state lives in SSZ linear memory (WASM) with typed-array or proxy views; amounts are fixed-point \code{bigint} or \code{u128}, never \code{number}; collections serialise in a specified canonical order. A host object model is a view, never the source of truth.
\end{enumerate}

\subsection{Developer experience}
\begin{lstlisting}[language=TypeScript,caption={Illustrative usage (names are proposals). Identifiers are placeholders supplied by the application, not fixed values.}]
const actor = await createSovereignActor({
  schema: erpSchema,                        // generated, end-to-end types
  storage: new IndexedDbStorage(appName),   // local first
  sync: new IrohSyncAdapter({ lane: "urgent" }),
});

const tx = await actor.reconcile((draft) => {
  // Bilateral appointment tuple in Slot 1 and a typed row anchored in Slot 4,
  // applied as one atomic virtual batch.
  draft.slot1.zanzibar.grant({ subject: anchorDaoDid, relation: "appointed_anchor", object: draft.did });
  draft.slot4.tables.invoices.insert({ id: invoiceId, debtor: debtorDid, amount: amountCents /* bigint */ });
}, (s) => ui.show(s.phase));
\end{lstlisting}

\subsection{Why the boundary matters}
Early Ethereum let implementation quirks (storage layout, provider-injection conventions, JSON-RPC assumptions) harden into de facto rules. Keeping the normative machine in the ledger paper and this SDK in the Stack paper is the guard against repeating that: the SDK may change freely so long as it passes the conformance vectors.

\section{Open decisions before migration}
\begin{itemize}
  \item Which build and attestation claims can be backed by a reproducibility report today, and which remain specification.
  \item Whether the Stack paper should also govern the paper build pipeline itself (this document set is built on the remote host and checksummed).
\end{itemize}

\end{document}
